\begingroup\small\setlength{\tabcolsep}{3pt}\renewcommand{\arraystretch}{1.15}
\begin{longtable}{>{\raggedright\arraybackslash}p{.14\textwidth}>{\raggedright\arraybackslash}p{.38\textwidth}>{\raggedright\arraybackslash}p{.39\textwidth}}
\caption{Prior-work actions and comparison boundaries. Names link to the primary sources; versions and locators are recorded in T3/data.json.}\label{tab:prior}\\
\toprule Method & State and action & Boundary of this comparison \\ \midrule\endfirsthead
\toprule Method & State and action & Boundary of this comparison \\ \midrule\endhead
\href{https://arxiv.org/html/2502.07115v5}{MC-SF} & Keep admitted jobs running; select waiting jobs using predicted upper lengths and future memory feasibility. & Targets total completion latency in a non-preemptive round model. Our FIFO peak adapter omits shortest-first subset selection; it is not full MC-SF. \\
\href{https://arxiv.org/abs/2504.11320v4}{WAIT} & Gate admission with fluid-guided thresholds; Nested WAIT regulates decode segments and adds a memory safety buffer. & Known-length and unknown-length variants target stability and latency. Memory headroom is prior art; our measurements concern native offload recovery. \\
\href{https://arxiv.org/abs/2404.16283}{Andes} & Prioritize and preempt at token boundaries using expected user-experience gain and GPU resource demand. & Targets smooth streaming and prompt first output. Our fixed maxgap/stall metrics do not reproduce its reading-aware utility. \\
\href{https://arxiv.org/abs/2305.05920}{FastServe} & Use skip-join multilevel feedback queues initialized from prompt length, with token-level preemption and proactive host/GPU state movement. & Targets job latency through priority and execution quanta. HORIZON retains native quanta and changes admission/recovery/victim decisions. \\
\href{https://arxiv.org/html/2411.18424v1}{FastSwitch} & Group KV blocks, adapt synchronous/asynchronous swaps, and reuse valid CPU KV copies to reduce switching cost. & Changes layout and transfer execution for fairness-aware serving. Our fixed native transfer path cannot reproduce its system-level benefit. \\
\href{https://arxiv.org/html/2606.18431v1}{UniBoost} & Unify prefill/decode priorities from arrival and attained work; quantize revisions and protect a dispatched request until its next service threshold. & MemGuard already provides minimum-run protection. Our step-based lease and peer-growth gate are a different implementation, not proof of novelty or superiority. \\
\href{https://arxiv.org/html/2601.20309v1}{SuperInfer} & Rotate requests by SLO-derived virtual lag under a transfer budget; overlap full-duplex KV movement with computation. & Designed for GH200 HBM/Grace memory and C2C transfers. Its proactive rotary scheduling and transfer engine are not reproduced by our single-GPU adapter. \\
\href{https://www.usenix.org/conference/osdi24/presentation/sun-biao}{Llumnix} & Reschedule running requests across model instances using live migration for load balancing, isolation, and priority. & Requires multiple instances and migration. Our experiment has one serving GPU; its remaining loss does not test distributed rescheduling. \\
\href{https://arxiv.org/abs/2309.06180}{vLLM} & Allocate KV in pages and share blocks to reduce fragmentation and duplicate storage. & The original paper motivates memory management. Native victim/LOAD/ACK behavior here is established from pinned 0.26 source, not attributed to the 2023 paper. \\
\bottomrule\end{longtable}
\endgroup
